\chapter{The Model, and the Object of Study}\label{ch:model}
\label{ch:background}\label{ch:graphs}\label{ch:graphical}\label{ch:development}

This chapter fixes what the thesis computes and the language it computes in.
A Markov chain generates a clean trajectory; an Ornstein--Uhlenbeck channel
corrupts every frame independently; the object of study is the joint score of
the corrupted sequence. Two justifications come before any analysis. The first
is that the quantity worth studying is the joint score of a whole trajectory
and not a per-frame score; the second is that the stationary AR(1) chain is the
right model in which to study it. Both are settled here. The chapter closes with the representation the rest of
the thesis reasons in: a distribution written as a product of local factors,
and the recursion that returns its marginals exactly when those factors form a
tree.

\section{Markov sequences}\label{sec:stationarity}\label{sec:markov}

Two notions of time homogeneity are used constantly below and have to be kept
apart \citep{brockwell1991time}. A process $(X_k)$ is \emph{strictly
stationary} if every finite-dimensional law is shift-invariant, and \emph{weakly
stationary} if its mean is constant and its autocovariance depends only on the
lag, $\operatorname{Cov}(X_k, X_{k+d}) = \gamma(d)$. Strict implies weak
whenever second moments exist, and for Gaussian processes the two coincide. The
distinction earns its place because Chapter~\ref{ch:laplace} replaces the
Gaussian innovation with a Laplace one, at which point they separate and only
the strict notion survives the construction. Weak stationarity also has the
matrix consequence the research chapters exploit: the covariance of $L$
consecutive observations, $(\Sigma)_{ij} = \gamma(i-j)$, is constant along
diagonals. Wherever a covariance in this thesis is Toeplitz, stationarity is the
reason.

\begin{definition}[Markov chain]
A sequence $(X_k)_{k \geq 0}$ on a state space $\mathcal{X}$ is a \emph{Markov
chain} if for every $k$ and every measurable $A \subseteq \mathcal{X}$,
\begin{equation}
  \Prob\big(X_{k+1} \in A \,\big|\, X_k, X_{k-1}, \dots, X_0\big)
  \;=\; \Prob\big(X_{k+1} \in A \,\big|\, X_k\big) .
  \label{eq:markov-property}
\end{equation}
The conditional law is encoded by a \emph{transition kernel}
$\kernel(x' \,|\, x)$, with $\int \kernel(x'|x)\,\dd x' = 1$ for every $x$.
\end{definition}

\begin{example}[The Gaussian AR(1) chain]\label{ex:ar1}
The running example throughout is the scalar Gaussian autoregression of order
one,
\begin{equation}
  a_{k+1} \;=\; \alpha\, a_k + \eta_k ,
  \qquad
  \eta_k \stackrel{\text{iid}}{\sim} \Nd(0, \ivar),
  \qquad k = 0, \dots, L-2 ,
  \label{eq:ar1}
\end{equation}
with kernel $\kernel(a' \mid a) = \Nd(a'; \alpha a, \ivar)$. Taking
$\ivar = 1 - \alpha^2$ and $a_0 \sim \Nd(0,1)$ makes it stationary with unit
marginal variance and $\gamma(d) = \alpha^{|d|}$, so its covariance is Toeplitz.
Every chapter from Chapter~\ref{ch:gaussian} onwards concerns \eqref{eq:ar1},
one of its non-Gaussian variants, or what an estimator can recover of it
through noise.
\end{example}

The Markov property \eqref{eq:markov-property} is exactly a statement about
\emph{factorisation}: the joint law of the first $L$ states is
\begin{equation}
  p(x_0, x_1, \dots, x_{L-1})
  \;=\; p_0(x_0)\, \prod_{k=0}^{L-2} \kernel(x_{k+1} \,|\, x_k) .
  \label{eq:chain-factorisation}
\end{equation}
This is the most heavily used equation of the thesis. It says the joint
distribution of a Markov trajectory is a product of \emph{local} terms, each
touching two consecutive variables: the structure that makes the chain a tree
in the factor-graph sense of Chapter~\ref{ch:graphs}, that makes its precision
matrix tridiagonal in the Gaussian case of Chapter~\ref{ch:gaussian}, and that
makes belief propagation exact. Marginals evolve by integrating the kernel,
$p_{k+1}(x') = \int \kernel(x'|x)\, p_k(x)\,\dd x$, and a law $\pi$ with
$\pi = \kernel \pi$ is \emph{stationary}: a chain started in $\pi$ remains in
it, and its trajectory is then strictly stationary.

One change of perspective sets up the research problem. Markov chains are
classically studied as \emph{algorithms}, where the stationary law is the target
and the trajectory a means to reach it. From Chapter~\ref{ch:gaussian} onwards
the chain \eqref{eq:ar1} plays the opposite role: it is the
\emph{data-generating process}, one trajectory $(a_0, \dots, a_{L-1})$ is one
data point, and the object of interest is the law of the whole trajectory after
every coordinate has been independently corrupted. The factorisation
\eqref{eq:chain-factorisation} then stops being a computational convenience and
becomes the structure to be recovered.

\section{The Ornstein--Uhlenbeck corruption channel}\label{sec:ou}

The noising process used by diffusion models, and by every result of this
thesis, is the \emph{Ornstein--Uhlenbeck process}
\citep{uhlenbeck1930theory}: Brownian motion pulled back to the origin by a
linear spring,
\begin{equation}
  \dd X_t \;=\; -X_t\,\dd t \;+\; \sqrt{2}\;\dd W_t ,
  \qquad X_0 = a .
  \label{eq:ou-sde}
\end{equation}
The drift and diffusion constants are the variance-preserving normalisation of
the diffusion literature \citep{song2021scorebased}; any other choice rescales
time. The process is the unique Gaussian, Markov, stationary process in
continuous time, and it is exactly solvable. Solving \eqref{eq:ou-sde} by
integrating factor, the derivation is routine and is carried out in
\compendium{}, gives a Gaussian transition kernel,
$p(x_t \mid X_0 = a) = \Nd(x_t;\, a e^{-t},\, \Deltat)$ with
$\Deltat = 1 - e^{-2t}$. Corrupting a data point $a$ for a time $t$ therefore
produces
\begin{equation}
  x \;=\; a\, e^{-t} \;+\; \sqrt{\Deltat}\;\varepsilon,
  \qquad \varepsilon \sim \Nd(0, 1):
  \label{eq:ou-channel}
\end{equation}
the signal is \emph{contracted} by $e^{-t}$ while isotropic noise of variance
$\Deltat$ fills in. At $t = 0$ the channel is the identity; as $t \to \infty$,
$x \sim \Nd(0,1)$ regardless of $a$. The signal-to-noise ratio
$e^{-2t}/\Deltat$ sweeps continuously from $\infty$ to $0$ and diverges as
$1/2t$ for small $t$, a factor that reappears throughout the thesis as an error
amplifier. The channel interpolates between the data distribution and a
standard Gaussian, which is what a diffusion model needs.

Two structural remarks are used repeatedly later. First, the kernel is Gaussian
\emph{in $a$ as well as in $x$}: read as a function of the clean variable,
$\Nd(x; a e^{-t}, \Deltat) \propto
\exp[-(e^{-2t}/2\Deltat)(a - x e^{t})^2]$, a Gaussian likelihood factor of
precision $e^{-2t}/\Deltat$ centred at the deconvolved observation $e^t x$.
This is what makes the posteriors of Chapter~\ref{ch:gaussian} jointly Gaussian
and the belief-propagation messages closable in a finite-dimensional family.
Second, if $a \sim \Nd(\mu, \Sigma_0)$ is pushed through \eqref{eq:ou-channel}
coordinate-wise with \emph{independent} noises, the output is Gaussian with
\begin{equation}
  \mu_t = e^{-t}\mu,
  \qquad
  \Sigt \;=\; e^{-2t}\, \Sigma_0 \;+\; \Deltat\, I ,
  \label{eq:ou-cov-map}
\end{equation}
since means and covariances transform linearly and the added noise is white.
Equation \eqref{eq:ou-cov-map}, shrink the covariance, add a multiple of the
identity, is the entire forward process of this thesis in one line.


\section{How the object of study was arrived at}\label{sec:dev-history}

The estimand of this thesis was not chosen at the outset. It was reached by
getting it wrong first, and the sequence is worth recording, both because it
explains the shape of what follows and because the correction is the single
most consequential step in the project.

\paragraph{The propagator hypothesis.} The work began from a structural guess
about sequences. If a trajectory is generated by a fixed dynamical rule, the score at one frame
ought to be obtainable from the score at the previous one by a fixed operator.
Writing $S_{:,u}$ for the block of the score at internal time $u$, the
conjecture was that some \emph{propagator} $P_t$ satisfies
\begin{equation}
  S_{:,u}(x) \;\approx\; P_t\big(S_{:,u-1}(x)\big).
  \label{eq:dev-propagator}
\end{equation}
If such an operator existed and were cheap, generation over a long sequence
would cost one solve plus $L-1$ applications of $P_t$. The programme adopted to
test \eqref{eq:dev-propagator} was deliberately austere: exactly solvable
models only, closed forms wherever possible, and no neural networks anywhere,
so that any structure found could be attributed to the model and not to a
fitting procedure.

\paragraph{Six toy models, and a mistake they shared.} A sequence of solvable
objects was built to probe \eqref{eq:dev-propagator}: a static Gaussian
mixture under the channel; a two-frame joint law; a two-dimensional rotating
process $a_{u+1} = R_{\Delta\theta}\, a_u + \xi_u$; a billiard in a box; a
scalar AR(1) chain examined on the whole $(u,t)$ plane; and a diffusion on the
circle. The rotating model and the circle are the direct ancestors of
Chapter~\ref{ch:ring}, and the AR(1) study is the ancestor of
Chapters~\ref{ch:gaussian} and~\ref{ch:gaussian-bp}.

All but one of them computed the wrong object. What was evaluated was the
\emph{per-frame marginal} score
\begin{equation}
  s(x_u, u, t) \;=\; \nabla_{x_u} \log p_t(x_u \mid u),
  \label{eq:dev-marginal-score}
\end{equation}
which integrates out every other frame. Doing so is not a mild approximation.
It is equivalent to asserting that the clean law factorises as
$P_0(a_0,\dots,a_{L-1}) = \prod_u p_0(a_u \mid u)$, which is false for any
correlated process and, in particular, false for every model in the list.
The quantity being studied had the frames' dependence removed from it before
the study began, so no propagator could have been found in it: there was
nothing left for $P_t$ to propagate.

\paragraph{The correction.} The object was corrected in a supervision meeting
of March 2026 to the \emph{joint} score of the whole trajectory,
$S(x,t) = \nabla_{(x_0,\dots,x_{L-1})} \log P_t(x_0,\dots,x_{L-1})$, whose
$k$-th block is the posterior expectation
\begin{equation}
  S_k(x,t) \;=\; \frac{e^{-t}\,\E\big[a_k \,\big|\, x_0,\dots,x_{L-1}\big] - x_k}{\Deltat},
  \label{eq:dev-joint-block}
\end{equation}
taken over the full posterior of the clean trajectory given all $L$
observations. Equation~\eqref{eq:dev-joint-block} is
\eqref{eq:tweedie-joint} of Chapter~\ref{ch:statmech}, and the conditioning on
every coordinate is the whole of the difference between it and
\eqref{eq:dev-marginal-score}. The remainder of this chapter makes the case for
it directly, without reference to how it was reached.

\paragraph{What the correction bought.} It made the stationary Gaussian AR(1)
chain the natural testbed rather than one model among six. Under
\eqref{eq:dev-joint-block} the posterior is Gaussian, so the conditional mean
is available from a Kalman smoother, the joint score is linear in the whole of
$x$, and the clean precision is tridiagonal. Each of these is a closed form
that can be checked against another closed form, which is what
Chapters~\ref{ch:gaussian} and~\ref{ch:gaussian-bp} do.

\paragraph{The propagator question, answered.} The conjecture that started the project was then settled rather than
abandoned, and the rest of the thesis is organised around the answer. A
propagator does exist for the Gaussian chain, and it is exact: the
Rauch--Tung--Striebel smoother reproduces the precision-matrix score to machine
precision, which Chapter~\ref{ch:gaussian-bp} proves and measures. What is
\emph{not} exact is the naive reading of \eqref{eq:dev-propagator} in which
$P_t$ acts as a shift, one row of the precision matrix mapping to the next.
That form holds only asymptotically: the precision is Toeplitz in the bulk but
not at the boundaries, and its deviation from Toeplitz falls with both chain
length and diffusion time, from a few per cent on short chains at low noise to
a small fraction of one per cent on long chains at high noise. A related
negative result closes the cheapest version of the idea: the joint law admits
no exact factorisation into independent one-dimensional bonds for any $L \geq
2$, because the cross term in the two-frame characteristic function is
$-\alpha\Deltat$, which vanishes only when $\alpha = 0$ or $t = 0$.

So the propagator is a smoother, not a shift. That is a less convenient answer
than the one hoped for, and a considerably more useful one, because a smoother
is exactly the object that generalises: it survives the loss of Gaussianity as
belief propagation (Chapter~\ref{ch:laplace}), and it survives the loss of a
known kernel as EM (Chapter~\ref{ch:em-method}).

\section{Why the joint score}\label{sec:dev-correction}

Standard diffusion models generate \emph{static} objects; the data of
interest here are \emph{dynamic}, a data point being a sequence of frames
indexed by an internal time $u = 0, 1, \dots, L-1$, a video, a rotating
shape, a stochastic trajectory. Two times therefore appear from the outset,
and keeping them apart is a permanent discipline of this thesis: the
\emph{internal} time $u$ of the sequence, and the \emph{diffusion} time $t$
of the noising channel, which acts independently on every frame.

It is tempting to treat such a trajectory as $L$ separate samples, computing
each frame's noisy marginal score independently. In the Gaussian model that
carries most of this thesis, that is not merely lossy, it is empty.

\begin{proposition}[One-frame marginals in the stationary Gaussian
chain]\label{prop:dev-marginal}
Let the clean chain be stationary Gaussian AR(1) with unit marginal variance,
put through the variance-preserving channel
$x = e^{-t} a + \sqrt{1 - e^{-2t}}\, z$, with $z \sim \Nd(0,1)$ independent of
$a$. Then every noisy frame is marginally $\Nd(0,1)$ at every diffusion time,
its per-frame marginal score is exactly $-x_k$, and the collection of one-frame
marginals carries zero Fisher information about the coupling $\alpha$.
\end{proposition}

\noindent This is immediate: variance preservation gives
$\operatorname{Var}(x) = e^{-2t}\operatorname{Var}(a) + (1 - e^{-2t}) = 1$, and a Gaussian $a$ makes $x$
Gaussian, so the marginal law is $\Nd(0,1)$ whatever $\alpha$ is and cannot
depend on it.

\begin{remark}[The same statement is false off the Gaussian
model]\label{rmk:dev-nongaussian}
Variance preservation alone gives only $\operatorname{Var}(x) = 1$; it does not give
Gaussianity, and the distinction matters because the non-Gaussian chains of
Chapters~\ref{ch:laplace}--\ref{ch:em-results} are the ones the rest of this
thesis measures. Writing $\varphi$ for characteristic functions,
\begin{equation}
  \varphi_{x}(u) \;=\; \varphi_a\!\big(e^{-t}u\big)\,
  \exp\!\Big[-\tfrac{1 - e^{-2t}}{2}\,u^2\Big],
  \label{eq:dev-cf}
\end{equation}
so if $x$ were exactly $\Nd(0,1)$ at some finite $t$, dividing by the Gaussian
factor would force
$\varphi_a(e^{-t}u) = \exp[-\tfrac{1}{2}e^{-2t}u^2]$, that is, $a$ Gaussian.
A genuinely non-Gaussian frame therefore stays non-Gaussian under any finite
amount of Gaussian corruption, and its one-frame score stays nonlinear;
Section~\ref{sec:l-K1} computes that score for a Laplace prior and shows it
approaching $-x$ only as $t \to \infty$.
\end{remark}

\noindent So the correct general statement is narrower than
``per-frame is empty'', and it is about \emph{dependence} and not about
information as such. One-frame marginals can see the innovation law, and can
see a parameter that changes the marginals; what no collection of them can
recover is how consecutive frames are coupled, since coupling is a property of
the joint law and of no marginal of it. In the Gaussian case those two
statements collapse into one, because the marginals are $\alpha$-free by
Proposition~\ref{prop:dev-marginal}; off it they come apart, and
Chapter~\ref{ch:ring} proves the sharper version for one specific parameter, that the rotation rate of the ring model has exactly zero Fisher information in
every one-frame marginal, under that model's rotational symmetry.

A sequence is \emph{one} sample from a joint law, and the quantity that
drives the reverse dynamics is the joint score of the whole noisy trajectory,
the gradient of $\log P_t(x_0, \dots, x_{L-1})$, whose $k$-th component is a
conditional expectation over the \emph{full} posterior of the clean frames,
as in \eqref{eq:tweedie-joint}. That is the object in which the temporal
structure lives. Chapter~\ref{ch:gaussian} exhibits the contrast in closed form
at $L = 2$.

This also settles the shape of the structural question that opened
Section~\ref{sec:dev-history}. The propagator \eqref{eq:dev-propagator} was
hoped for so that a learning system could hold one block and one map instead of
$L$ blocks. What the joint view shows is that the structure worth studying is
not a map between per-frame scores at all, but the coupling pattern of the
joint score: a statement about a matrix in the Gaussian case, and about a
posterior in general. That is the object the rest of this thesis analyses.
Figure~\ref{fig:toymodel} shows it in the smallest case where it exists: the
diagonal tilt of the joint score field is the whole of the coupling, and it is
exactly what a per-frame treatment throws away.

\begin{figure}[t]\centering
\includegraphics[width=\linewidth]{fig_toymodel_score.pdf}
\caption[The coupling made visible at $L=2$]{The coupling made visible, in the smallest case where it exists:
$x_1 = x_0 + c + \eta$ with a trimodal Gaussian-mixture prior on $x_0$, drift
$c = 1$, innovation standard deviation $\sigma = 0.45$, after diffusion time
$t = 0.15$. Left: the noised joint density $p_t(x_0, x_1)$. Right: the joint
score field $\nabla \log p_t$ (arrows) over the same density (contours). Both
panels are computed from the closed-form mixture expressions on a grid. The
arrows point towards the modes and tilt along the diagonal, and that tilt is
the entire effect of the coupling between frames, it is exactly what any
per-frame treatment discards. Note that this prior is \emph{not} Gaussian, so
by Remark~\ref{rmk:dev-nongaussian} its noised one-frame marginals are Gaussian
mixtures with nonlinear scores rather than $\Nd(0,1)$ with score $-x_k$; what
those marginals cannot see is the coupling, which is the point the figure
makes.}
\label{fig:toymodel}
\end{figure}

\section{The minimal model}\label{sec:dev-minimal}

Several dynamic objects were candidates: rotations, coupled phases, bounded
trajectories with reflections, and curved data manifolds. The stationary
Gaussian AR(1) chain under coordinatewise OU noise is the minimal one in which
the joint score is nontrivial yet exactly computable. one coupling
parameter $\alpha$ for the internal dynamics against one noise parameter $t$
for the corruption, linear algebra available at every step, and every
structural question (coupling range, locality, message passing) posable in
closed form.

The requirement doing the work here is \emph{exact computability}. The whole
programme of this thesis is to measure what approximations cost, Gaussian
message closure in Chapter~\ref{ch:laplace}, an estimated transition kernel
in Chapters~\ref{ch:em-method} and~\ref{ch:em-results}, a finite receptive
field in Chapter~\ref{ch:gaussian}, and every one of those measurements
needs a reference that is known rather than approximated. A model in which
the answer is available in closed form is not a simplification of the
question; it is what makes the question answerable at all.

The nonlinear models are set aside at that point deliberately, as targets for
after the solvable case is understood. Chapter~\ref{ch:ring} takes up the
first of them.


\section{Markov random fields and factor graphs}\label{sec:mrf}
\label{sec:factor-graphs}

Graphical models assign structure to distributions the way a precision matrix
assigns it to a Gaussian, and for arbitrary variables. A collection
$x = (x_i)_{i \in V}$ on an undirected graph $G$ is a \emph{Markov random field}
if each variable is conditionally independent of the rest given its neighbours,
$p(x_i \mid x_{V \setminus \{i\}}) = p(x_i \mid x_{\partial i})$. For a
strictly positive density the Hammersley--Clifford theorem
\citep{besag1974spatial} makes that equivalent to factorisation over the cliques
of $G$,
\begin{equation}
  p(x) \;=\; \frac{1}{Z} \prod_{C \in \mathrm{cliques}(G)} \psi_C(x_C)
  \;=\; \frac{1}{Z}\, \exp\Big[ -\sum_C E_C(x_C) \Big] ,
  \label{eq:hammersley-clifford}
\end{equation}
so conditional independence, graphical locality and additive local energies are
one thing seen three ways. For Gaussian fields it becomes matrix-concrete: the
energy is $\tfrac12 a^\T J a - h^\T a$, and the graph is exactly the sparsity
pattern of $J$. A Gaussian Markov chain therefore has a tridiagonal precision,
which Chapter~\ref{ch:gaussian} derives and uses throughout.

One consequence carries the whole thesis. If each variable is observed through
its own independent channel, the likelihood factors touch one variable each and
no clique is enlarged: \emph{the posterior of an MRF under per-node observations
is an MRF on the same graph}. Observation changes the potentials, never the
graph.

An MRF records which variables interact but not how the joint factorises: a
triangle of pairwise terms and one triple term have the same graph. The
\emph{factor graph} makes the factorisation itself graphical
\citep{kschischang2001factor}. Given
$p(x) = \frac{1}{Z}\prod_a \psi_a(x_{\partial a})$, it is the bipartite graph
with a variable node per $x_i$, a factor node per $\psi_a$, and an edge whenever
$x_i \in x_{\partial a}$. A connected factor graph with one fewer edge than
nodes is a \emph{tree}, and trees are the tractable case: on them marginals and
partition functions are exact under the recursion of the next section, whatever
the alphabet. The factor graph of a Markov chain observed through per-node
noise is a tree of a particularly simple shape, a spine of variables and
transition factors with one evidence leaf per variable, as
Chapter~\ref{ch:gaussian} verifies by counting.

\section{The sum--product algorithm}\label{sec:sum-product}

On a factor graph, sum--product belief propagation (BP) passes messages
along every edge \citep{pearl1988probabilistic, kschischang2001factor,
mezard2009information}. Variable-to-factor messages multiply the
incoming factor messages from all \emph{other} neighbours, and
factor-to-variable messages integrate the factor against the incoming
variable messages:
\begin{equation}
  \nu_{i \to a}(x_i) \propto
    \prod_{b \in \partial i \setminus a} \hat\nu_{b \to i}(x_i),
  \qquad
  \hat\nu_{a \to i}(x_i) \propto
    \int \psi_a(x_{\partial a})
    \prod_{j \in \partial a \setminus i} \nu_{j \to a}(x_j)\;
    \dd x_{\partial a \setminus i},
  \label{eq:bp-rules}
\end{equation}
with the belief at node $i$ proportional to
$\prod_{a \in \partial i} \hat\nu_{a \to i}(x_i)$. The ``exclude the
recipient'' rule is the heart of the algorithm: the message $i \to a$
summarises everything node $i$ knows \emph{except} what it learned from
$a$, so that no piece of evidence is ever counted twice along a path.

\begin{theorem}[Tree exactness of sum--product]\label{thm:tree-exact}
On a tree-shaped factor graph, one inward and one outward pass of
\eqref{eq:bp-rules} compute the exact node and factor marginals of $p$
\citep[Theorem~14.1]{mezard2009information}. The resulting beliefs are
unique; the individual unnormalised messages are defined only up to positive
multiplicative constants.
\end{theorem}

\emph{Proof.} Each edge of a tree splits the graph into two disjoint halves, so
the message crossing it summarises exactly the half behind it and no evidence
reaches a node twice. The recursion terminates with every marginal having
absorbed each factor once. \hfill$\square$

 On graphs with cycles the same updates,
iterated, define \emph{loopy} BP, which is extraordinarily useful but
approximate \citep{weiss2001correctness, malioutov2006walk}; none of that
is needed in this thesis, because our graphs are trees.

\begin{figure}[t]\centering
\begin{tikzpicture}[
  latent/.style={circle, draw, minimum size=7mm, inner sep=1pt},
  obs/.style={circle, draw, fill=black!10, minimum size=7mm, inner sep=1pt},
  fac/.style={rectangle, draw, fill=black, minimum size=2.2mm, inner sep=0pt},
  node distance=9mm, scale=0.92, every node/.style={transform shape}]
\foreach \i in {1,2,3,4}{
  \node[latent] (a\i) at (2.1*\i, 0) {$a_{\i}$};
  \node[fac] (f\i) at (2.1*\i, -1.0) {};
  \node[obs] (x\i) at (2.1*\i, -2.0) {$x_{\i}$};
  \draw (a\i) -- (f\i) -- (x\i);
}
\foreach \i/\j in {1/2,2/3,3/4}{
  \node[fac] (w\i) at (2.1*\i+1.05, 0) {};
  \draw (a\i) -- (w\i) -- (a\j);
}
\node[above=1mm of w1] {\small $\kernel(a_{k}\mid a_{k-1})$};
\node[right=2mm of f4] {\small $\like(x_k\mid a_k)$};
\draw[-{Stealth}, thick] (1.15,1.35) -- (2.55,1.35) node[midway,above] {\small $\lmsg$};
\draw[-{Stealth}, thick] (9.35,1.35) -- (7.95,1.35) node[midway,above] {\small $\rmsg$};
\end{tikzpicture}
\caption[Posterior factor graph of a noised Markov chain]{The posterior factor
graph of a Markov chain observed through the coordinatewise channel. Open
circles are latent variables, shaded circles observations, squares factors. Each
transition factor $\kernel(a_k \mid a_{k-1})$ joins two adjacent latent
variables; each observation enters through a \emph{unary} factor attached to a
single latent variable, because the channel corrupts every coordinate
independently. Conditioning therefore reweights node potentials without creating
any edge among the $a_k$, so the posterior graph is still a chain and
Theorem~\ref{thm:tree-exact} applies to it. Arrows show the two message sweeps
of \eqref{eq:bp-rules}. The construction fails as soon as the forward noise is
correlated across coordinates: the likelihood then touches several $a_k$ at
once and the graph acquires edges.}
\label{fig:posterior-factor-graph}
\end{figure}

\paragraph{The representation problem.} For \emph{discrete} variables
each message is a finite probability vector and each update in
\eqref{eq:bp-rules} is finite arithmetic. For \emph{continuous} variables
each message is a probability \emph{density}, an infinite-dimensional
object, and the update integrals must be carried out in function space.
Tree exactness (a property of the graph) says nothing about
representability of the messages (a property of the model's function
class); keeping the two apart is essential, and the distinction is where
all the difficulty of the non-Gaussian research chapter lives. There are
three ways to make continuous BP computable: restrict to models whose
messages close in a finite-dimensional family (the Gaussian route, next
section); represent messages numerically on a grid (the route our
non-Gaussian reference solver takes, Chapter~\ref{ch:laplace}); or
project each message onto a tractable family and accept the error (the
route whose cost this thesis measures).

\section{The Gaussian toolkit}\label{app:gauss}\label{sec:gauss-toolkit}

Five standard facts about Gaussian densities, collected here because
Chapters~\ref{ch:gaussian} to~\ref{ch:laplace} use them repeatedly and each
use is cleaner for pointing at a numbered identity than for restating it.
Nothing in this section is original.

This appendix collects, with proofs, the handful of Gaussian identities on
which Chapter~\ref{ch:gaussian} relies. Notation: a Gaussian
density in \emph{moment} parameters is $\Nd(a; m, v)$; in \emph{natural}
parameters it is
$\exp(-\tfrac{\lambda}{2} a^2 + \theta a + \mathrm{const})$ with
$\lambda = 1/v$, $\theta = m/v$.

\subsection{Quadratic forms and Gaussian normalisation}

For $J \succ 0$ (symmetric positive definite) and any $h$,
\begin{equation}
  \int_{\R^L} \exp\!\Big( -\tfrac{1}{2} a^\T J a + h^\T a \Big)\, \dd a
  = \frac{(2\pi)^{L/2}}{\sqrt{\det J}}
    \exp\!\Big( \tfrac{1}{2} h^\T J^{-1} h \Big),
  \label{eq:gauss-int}
\end{equation}
by completing the square, $-\tfrac{1}{2}a^\T J a + h^\T a =
-\tfrac{1}{2}(a - J^{-1}h)^\T J (a - J^{-1}h) + \tfrac{1}{2} h^\T J^{-1}h$,
and translating. Consequently a density
$p(a) \propto \exp(-\tfrac{1}{2} a^\T J a + h^\T a)$ is the Gaussian with
mean $m = J^{-1} h$ and covariance $J^{-1}$: \emph{reading off $(J, h)$
from a log-density is the fastest route to posterior parameters}, used in
Proposition~\ref{prop:g-posterior}.

\subsection{Products of Gaussians in the same variable}

\begin{equation}
  \Nd(a; m_1, v_1)\, \Nd(a; m_2, v_2)
  \;\propto\; \Nd\!\big(a;\, m, v\big),
  \qquad
  \frac{1}{v} = \frac{1}{v_1} + \frac{1}{v_2},
  \qquad
  \frac{m}{v} = \frac{m_1}{v_1} + \frac{m_2}{v_2}.
  \label{eq:gauss-product}
\end{equation}
\emph{Proof.} Add the exponents; the sum of two quadratics in $a$ is a
quadratic in $a$ whose $a^2$-coefficient is
$-\tfrac{1}{2}(1/v_1 + 1/v_2)$ and whose linear coefficient is
$m_1/v_1 + m_2/v_2$; apply \eqref{eq:gauss-int} in one dimension. In
natural parameters the rule is addition:
$(\lambda, \theta) = (\lambda_1 + \lambda_2, \theta_1 + \theta_2)$,
which is why the BP updates of Chapter~\ref{ch:gaussian} are two
additions per message. $\square$

\subsection{Affine propagation (marginalisation)}

If $a' = \alpha a + \eta$ with $\eta \sim \Nd(0, \sigma^2)$ independent of
$a \sim \Nd(m, v)$, then
\begin{equation}
  \int \Nd(a'; \alpha a, \sigma^2)\, \Nd(a; m, v)\, \dd a
  = \Nd\!\big(a';\, \alpha m,\, \alpha^2 v + \sigma^2\big).
  \label{eq:gauss-predict}
\end{equation}
\emph{Proof.} $a'$ is a sum of independent Gaussians, hence Gaussian; its
mean is $\alpha m$ and its variance $\alpha^2 v + \sigma^2$ by linearity
of expectation and independence. Read in the reverse direction (solving
for $a$ given a Gaussian summary of $a'$), the same computation gives
$\Nd\big(a;\, m'/\alpha,\, (v' + \sigma^2)/\alpha^2\big)$, the backward
half-step of the BP sweep. $\square$

\subsection{Gaussian conditioning}

For a jointly Gaussian pair with precision blocks
$\begin{psmallmatrix} J_{11} & J_{12} \\ J_{12}^\T & J_{22}
\end{psmallmatrix}$, the conditional of block $1$ given block $2$ is
Gaussian with precision $J_{11}$ (conditioning \emph{reads off} the
precision block) and mean shifted by $-J_{11}^{-1} J_{12}\, a_2$;
marginalisation, dually, is trivial in covariance parameters. This
precision/covariance duality (conditioning is local in $J$,
marginalising is local in $\Sigma$) is the algebraic engine behind the
cavity recursion recorded in \compendium{}: integrating a neighbour out of a
tridiagonal precision produces the rank-one correction
$-\beta^2/\lambda$ on the adjacent diagonal entry, by the Schur
complement
\begin{equation}
  \begin{pmatrix} \lambda & \beta \\ \beta & J_d \end{pmatrix}
  \;\longrightarrow\;
  J_d - \frac{\beta^2}{\lambda}.
  \label{eq:gauss-schur}
\end{equation}

\subsection{Woodbury and the two routes to the score}

The Woodbury identity
$(A + UCV)^{-1} = A^{-1} - A^{-1}U (C^{-1} + V A^{-1} U)^{-1} V A^{-1}$,
specialised to $\Sigt = \Deltat I + e^{-2t} \Sigz$, connects the noisy
precision $\Qt$ and the posterior precision
$J = (e^{-2t}/\Deltat) I + \Qz$:
\begin{equation}
  \Qt
  = \frac{1}{\Deltat}
    \left( I - \frac{e^{-2t}}{\Deltat}\, J^{-1} \right),
  \label{eq:gauss-woodbury}
\end{equation}
which is exactly the operator form of the score--posterior identity: substituting the
posterior mean $m = J^{-1} h = J^{-1} (e^{-t}/\Deltat) x$ into
$S = (e^{-t} m - x)/\Deltat$ gives $S = -\Qt x$ identically. The
machine-precision agreement of the two score routes
(Section~\ref{sec:g-score}) is the numerical shadow of
\eqref{eq:gauss-woodbury}.

\section{Locality, and what this thesis can check}\label{sec:cnn-locality}

Section~\ref{sec:sm-unet} set out the reading of U-Nets as belief propagation,
and the cross-architecture evidence for it. One consequence of that reading
belongs here, because it is what makes the chain worth studying rather than a
convenient simplification.

Those results concern hierarchical data, where the graphical model is a
balanced tree and the depth an architecture needs is logarithmic in the
sequence length. A chain is a degenerate tree, of depth equal to its length and
width two, so the same argument predicts a depth linear in the correlation
length instead. That is a sharply different architectural prediction, on a model
where it can be checked.

This thesis works one step below the correspondence, on a model small enough
that the locality of the score is a computable quantity and not an
architectural hypothesis. For the Gaussian chain the exact score at frame $k$
depends on the observation at frame $k \pm d$ with a weight decaying
geometrically in $d$, at a rate $q(\alpha,t)$ available in closed form
(Chapter~\ref{ch:gaussian-bp}); the error of the radius-$r$ windowed estimator
is measured to decay at the same rate, which is a separate and weaker
statement. Comparing truncation of the inference against truncation of the
score matrix then quantifies, in one solvable case, how a local computation
should be organised. All of this describes the function a network would have to
learn and not the dynamics of learning it; the trained baselines of
Chapter~\ref{ch:em-results} supply the other half.
